%% ----------------------------------------------------------------------------
%%  §4.3  Legame tra limiti di funzioni e limiti di successioni
%% ----------------------------------------------------------------------------
\section{Legame tra limiti di funzioni e limiti di successioni}
\label{sec:legame-successioni}

In questo paragrafo dimostriamo il \cref{teo:eps-delta}, cioè l'equivalenza tra la
definizione di limite «per successioni» (\cref{def:limite-funzione}) e quella con le
disuguaglianze. Fissati $f:A\to\R$, un punto $x_{0}\in\R$ di accumulazione per $A$ e
$\ell\in\R$, si tratta di provare che le due relazioni
\begin{align}
  &\perogni x_{n}\to x_{0},\quad x_{n}\in A\setminus\{x_{0}\}\ \perogni n\in\N
  \quad\implica\quad f(x_{n})\to\ell; \label{eq:ponte-succ}\\
  &\perogni\epsilon>0\ \esiste\delta>0 :\quad x\in A,\ \ 0<\abs{x-x_{0}}<\delta
  \quad\implica\quad \abs{f(x)-\ell}<\epsilon \label{eq:ponte-eps}
\end{align}
sono fra loro equivalenti.

\begin{schema}
\Ipotesi{$f:A\to\R$; \ $x_{0}\in\R$ di accumulazione per $A$; \ $\ell\in\R$.}
\Tesi{\eqref{eq:ponte-succ} \ $\sse$ \ \eqref{eq:ponte-eps}.}
\end{schema}

\begin{strategia}
Le due implicazioni hanno natura molto diversa. Che la \eqref{eq:ponte-eps} implichi la
\eqref{eq:ponte-succ} è quasi immediato: se una successione $x_{n}$ tende a $x_{0}$, i
suoi termini finiscono per stare nell'intorno di raggio $\delta$, e lì la
\eqref{eq:ponte-eps} controlla i valori di $f$. Il viceversa è più sottile, perché
dalla conoscenza di $f$ lungo le successioni dobbiamo ricavare un'informazione su
\emph{tutti} i punti di un intorno. Lo dimostriamo per assurdo: se la
\eqref{eq:ponte-eps} fosse falsa, esisterebbe una tolleranza $\epsilon_{0}$ violata in
ogni intorno di $x_{0}$, per quanto piccolo; scegliendo un punto «cattivo» in ciascuno
degli intorni di raggio $\frac1n$ costruiremo una successione che tende a $x_{0}$ lungo
la quale $f$ \emph{non} tende a $\ell$, contro la \eqref{eq:ponte-succ}.
\end{strategia}

\begin{dimostrazione}[Dimostrazione del \cref{teo:eps-delta}]
\emph{La \eqref{eq:ponte-eps} implica la \eqref{eq:ponte-succ}.} Consideriamo una
generica successione $x_{n}$ di punti di $A$, convergente a $x_{0}$, con $x_{n}\neq x_{0}$
per ogni $n\in\N$, e fissiamo $\epsilon>0$; sia $\delta>0$ il numero per cui vale la
\eqref{eq:ponte-eps}. Per la definizione di limite di successione (applicata con $\delta$
al posto di $\epsilon$) esiste un indice $\nu$ per cui $\abs{x_{n}-x_{0}}<\delta$ per ogni
$n>\nu$; inoltre, essendo $x_{n}\neq x_{0}$, in definitiva si ha
\begin{equation}\label{eq:ponte-passo1}
  x_{n}\in A, \qquad 0<\abs{x_{n}-x_{0}}<\delta \qquad\perogni n>\nu .
\end{equation}
Per l'ipotesi \eqref{eq:ponte-eps}, applicata con $x=x_{n}$, segue allora
\begin{equation}\label{eq:ponte-passo2}
  \abs{f(x_{n})-\ell}<\epsilon \qquad\perogni n>\nu,
\end{equation}
che, in base alla definizione di limite di successione, significa $f(x_{n})\to\ell$ per
$n\to+\infty$.

\emph{La \eqref{eq:ponte-succ} implica la \eqref{eq:ponte-eps}.} Procediamo per assurdo.
La \eqref{eq:ponte-eps} ha la forma «per ogni $\epsilon$ esiste $\delta$ tale che per
ogni $x$ vale una certa proprietà»; la sua negazione si ottiene scambiando i
quantificatori «per ogni» ed «esiste» e negando la proprietà finale. Contraddire la
\eqref{eq:ponte-eps} equivale dunque ad affermare che
\begin{equation}\label{eq:ponte-negazione}
\begin{split}
  \esiste\epsilon_{0}>0 :\ \perogni\delta>0\ \esiste x\in A :\quad &0<\abs{x-x_{0}}<\delta,\\
  &\abs{f(x)-\ell}\ge\epsilon_{0} .
\end{split}
\end{equation}
Poniamo $\delta=\frac1n$, con $n\in\N$, e indichiamo con $x=x_{n}$ un punto che compare
nella \eqref{eq:ponte-negazione} in corrispondenza di $\delta=\frac1n$. Otteniamo
\begin{equation}\label{eq:ponte-succ-cattiva}
  \esiste\epsilon_{0}>0 :\ \perogni n\in\N\ \esiste x_{n}\in A :\quad
  0<\abs{x_{n}-x_{0}}<\frac1n, \quad \abs{f(x_{n})-\ell}\ge\epsilon_{0} .
\end{equation}
Risulta in particolare
\begin{equation}\label{eq:ponte-carabinieri}
  x_{n}\neq x_{0}, \qquad x_{0}-\frac1n<x_{n}<x_{0}+\frac1n \qquad\perogni n\in\N;
\end{equation}
perciò $x_{n}\in A\setminus\{x_{0}\}$ per ogni $n$ e $x_{n}\to x_{0}$, per il teorema dei
carabinieri. Però $f(x_{n})$ non converge a $\ell$, perché la disuguaglianza
$\abs{f(x_{n})-\ell}\ge\epsilon_{0}$, valida per \emph{ogni} $n\in\N$, contrasta con la
definizione di limite di successione scritta con $\epsilon=\epsilon_{0}$ (non può
esistere alcun $\nu$ oltre il quale $\abs{f(x_{n})-\ell}<\epsilon_{0}$). Ciò contraddice
la \eqref{eq:ponte-succ}.
\end{dimostrazione}

\begin{osservazione}[Il caso dei punti e dei limiti infiniti]\label{oss:ponte-infinito}
La stessa dimostrazione vale, quasi parola per parola, quando $x_{0}$ oppure $\ell$, o
entrambi, sono infiniti, cioè prova l'equivalenza tra le definizioni
\eqref{eq:def-lim-piu-inf}--\eqref{eq:def-lim-inf-inf} e le corrispondenti formulazioni
della \cref{tab:intorni}. Il modo più rapido di vederlo è usare il linguaggio degli
intorni \eqref{eq:def-intorni}, osservando che le definizioni di limite di successione,
finito o infinito, del \cref{cap:successioni} si riassumono in un'unica frase:
\[
  x_{n}\to x_{0}\in\overline{\R} \quad\sse\quad
  \begin{gathered}[t]
    \text{per ogni intorno } U \text{ di } x_{0},\\
    x_{n}\in U \text{ definitivamente.}
  \end{gathered}
\]
Con questa frase, la prima parte della dimostrazione si ripete identica. Nella seconda,
la negazione della \eqref{eq:def-intorni} fornisce un intorno $V_{0}$ di $\ell$ tale che
ogni intorno $U$ di $x_{0}$ contiene un punto $x\in A\setminus\{x_{0}\}$ con
$f(x)\notin V_{0}$; scegliendo come $U$ l'intorno di raggio $\frac1n$ se $x_{0}\in\R$,
la semiretta $(n,+\infty)$ se $x_{0}=+\infty$ e la semiretta $(-\infty,-n)$ se
$x_{0}=-\infty$, si costruisce una successione $x_{n}\to x_{0}$ con $f(x_{n})\notin V_{0}$
per ogni $n$, che quindi non tende a $\ell$. Il lettore scriva per esteso il caso
$x_{0}=+\infty$, $\ell\in\R$ (\cref{es:ponte-infinito}).
\end{osservazione}

\subsection{Il limite è una proprietà locale}

Una conseguenza immediata della caratterizzazione con gli intorni è che il limite di $f$
per $x\to x_{0}$ dipende soltanto dai valori di $f$ \emph{vicino} a $x_{0}$: modificare
$f$ lontano da $x_{0}$ non cambia nulla.

\begin{lemma}[Località del limite]\label{lem:localita}
Siano $f:A\to\R$, $B\subseteq A$, e $x_{0}\in\overline{\R}$ di accumulazione per $B$; sia
$\ell\in\overline{\R}$.
\begin{enumerate}[label=\textup{(\roman*)}]
  \item Se $\lim_{x\to x_{0}}f(x)=\ell$, anche la restrizione di $f$ a $B$ tende a $\ell$
  per $x\to x_{0}$.
  \item Se esiste un intorno $U_{0}$ di $x_{0}$ tale che $U_{0}\cap A\subseteq B$, vale anche il
  viceversa: se la restrizione di $f$ a $B$ tende a $\ell$, allora
  $\lim_{x\to x_{0}}f(x)=\ell$.
\end{enumerate}
\end{lemma}

\begin{schema}
\Ipotesi{$B\subseteq A$; \ $x_{0}$ di accumulazione per $B$; \ (i) $f(x)\to\ell$; \
(ii) la restrizione di $f$ a $B$ tende a $\ell$, e $U_{0}\cap A\subseteq B$ per un
intorno $U_{0}$ di $x_{0}$.}
\Tesi{(i) la restrizione di $f$ a $B$ tende a $\ell$; \ (ii) $f(x)\to\ell$ per
$x\to x_{0}$.}
\end{schema}

\begin{strategia}
Si lavora con la definizione «per successioni». La (i) è immediata, perché le
successioni da controllare per la restrizione sono una parte di quelle da controllare
per $f$. Per la (ii) si osserva che una successione di punti di $A$ che tende a $x_{0}$
entra definitivamente nell'intorno $U_{0}$, e quindi in $B$: tolti i primi termini, che
non influiscono sul limite, è una successione di punti di $B$.
\end{strategia}

\begin{dimostrazione}
(i) Ogni successione di punti di $B\setminus\{x_{0}\}$ che tende a $x_{0}$ è anche una
successione di punti di $A\setminus\{x_{0}\}$: lungo di essa $f$ tende a $\ell$ per
ipotesi.

(ii) Sia $x_{n}\in A\setminus\{x_{0}\}$ con $x_{n}\to x_{0}$. Per la frase
dell'\cref{oss:ponte-infinito}, $x_{n}\in U_{0}$ definitivamente, diciamo per $n>\nu$;
quindi $x_{n}\in B$ per $n>\nu$. La successione traslata $x_{\nu+k}$, $k\in\N$, è
formata da punti di $B\setminus\{x_{0}\}$ e tende a $x_{0}$: per ipotesi
$f(x_{\nu+k})\to\ell$ per $k\to+\infty$. Ne segue $f(x_{n})\to\ell$, perché le due
successioni $f(x_{n})$ e $f(x_{\nu+k})$ differiscono soltanto per i primi $\nu$ termini:
se, ad esempio con $\ell\in\R$, si ha $\abs{f(x_{\nu+k})-\ell}<\epsilon$ per ogni $k>k_{0}$,
allora $\abs{f(x_{n})-\ell}<\epsilon$ per ogni $n>\nu+k_{0}$.
\end{dimostrazione}

La parte (ii) non vale senza l'ipotesi sull'intorno: la funzione $\frac{\abs x}{x}$,
ristretta a $B=(0,+\infty)$, tende a $1$ per $x\to0$, ma su tutto il suo dominio non ha
limite. Il limite destro e il limite sinistro sono proprio i limiti delle restrizioni a
$A\cap(x_{0},+\infty)$ e ad $A\cap(-\infty,x_{0})$; la proposizione seguente dice come
ricostruire da essi il limite.

\begin{proposizione}\label{prop:dx-sx}\index{limite!destro e sinistro}
Sia $x_{0}\in\R$ di accumulazione sia per $A\cap(x_{0},+\infty)$ sia per
$A\cap(-\infty,x_{0})$, e sia $\ell\in\overline{\R}$. Allora
\[
  \lim_{x\to x_{0}}f(x)=\ell \quad\sse\quad
  \lim_{x\to x_{0}^{+}}f(x)=\lim_{x\to x_{0}^{-}}f(x)=\ell .
\]
\end{proposizione}

\begin{schema}
\Ipotesi{$x_{0}\in\R$ di accumulazione per $A\cap(x_{0},+\infty)$ e per
$A\cap(-\infty,x_{0})$; \ $\ell\in\overline{\R}$.}
\Tesi{$f(x)\to\ell$ per $x\to x_{0}$ \ $\sse$ \ limite destro e limite sinistro esistono
e valgono entrambi $\ell$.}
\end{schema}

\begin{strategia}
Limite destro e sinistro sono limiti di restrizioni di $f$, e quindi la parte
«$\implica$» è la (i) del lemma di località. Per il viceversa si usano le disuguaglianze:
ciascuno dei due limiti unilaterali fornisce un $\delta$ dal proprio lato di $x_{0}$, e il
più piccolo dei due va bene da entrambi i lati.
\end{strategia}

\begin{dimostrazione}
L'implicazione «$\implica$» segue dal \cref{lem:localita}(i). Viceversa, supponiamo
$\ell\in\R$ (il caso infinito è analogo, con $M$ al posto di $\epsilon$) e fissiamo
$\epsilon>0$. Per la \eqref{eq:def-limite-dx} esiste $\delta_{1}>0$ tale che
$\abs{f(x)-\ell}<\epsilon$ per $x\in A$, $x_{0}<x<x_{0}+\delta_{1}$; per la
\eqref{eq:def-limite-sx} esiste $\delta_{2}>0$ tale che $\abs{f(x)-\ell}<\epsilon$ per
$x\in A$, $x_{0}-\delta_{2}<x<x_{0}$. Posto $\delta=\min\{\delta_{1},\delta_{2}\}$, per ogni
$x\in A$ con $0<\abs{x-x_{0}}<\delta$ si ha $x$ nel primo o nel secondo caso, e quindi
$\abs{f(x)-\ell}<\epsilon$. Per il \cref{teo:eps-delta}, $f(x)\to\ell$.
\end{dimostrazione}

\subsection{Come si dimostra che un limite non esiste}

La definizione «per successioni» è particolarmente efficace per dimostrare che un limite
\emph{non} esiste: basta esibire due successioni lungo le quali la funzione si comporta
in modo diverso.

\begin{corollario}[Criterio di non esistenza]\label{cor:non-esistenza}
Se esistono due successioni $x_{n}$, $x_{n}'$ di punti di $A\setminus\{x_{0}\}$, entrambe
tendenti a $x_{0}$, tali che $f(x_{n})\to\lambda$ e $f(x_{n}')\to\lambda'$ con
$\lambda\neq\lambda'$, allora il limite di $f(x)$ per $x\to x_{0}$ non esiste. Lo stesso
accade se esiste una sola successione $x_{n}\in A\setminus\{x_{0}\}$, $x_{n}\to x_{0}$,
lungo la quale $f(x_{n})$ non è regolare.
\end{corollario}

\begin{schema}
\Ipotesi{$x_{n},x_{n}'\in A\setminus\{x_{0}\}$, \ $x_{n}\to x_{0}$, $x_{n}'\to x_{0}$, \
$f(x_{n})\to\lambda$, $f(x_{n}')\to\lambda'$, $\lambda\neq\lambda'$ \ (oppure: $f(x_{n})$
non regolare).}
\Tesi{$\lim_{x\to x_{0}}f(x)$ non esiste.}
\end{schema}

\begin{strategia}
È la negazione della \cref{def:limite-funzione}: se il limite esistesse, la funzione
dovrebbe tendere a esso lungo \emph{tutte} le successioni ammissibili. Si ragiona per
assurdo e si conclude con l'unicità del limite di una successione.
\end{strategia}

\begin{dimostrazione}
Se esistesse $\lim_{x\to x_{0}}f(x)=\ell$, per la \cref{def:limite-funzione} sarebbe
$f(x_{n})\to\ell$ e $f(x_{n}')\to\ell$, quindi $\lambda=\ell=\lambda'$ per l'unicità del
limite di successione; nel secondo caso $f(x_{n})$ avrebbe limite $\ell$.
\end{dimostrazione}

\begin{esempio}\label{es:sen-cos-infinito}
I limiti
\begin{equation}\label{eq:sen-cos-infinito}
  \lim_{x\to+\infty}\sen x, \qquad \lim_{x\to+\infty}\cos x \qquad\text{non esistono.}
\end{equation}
Limitiamoci al primo dei due. Se esistesse il limite $\ell$, dovremmo avere che
$\sen x_{n}$ tende sempre allo stesso valore $\ell$ qualunque sia la successione
$x_{n}\to+\infty$. Ma ponendo $x_{n}=2\pi n$ e $x_{n}'=2\pi n+\frac\pi2$, entrambe
divergenti a $+\infty$, risulta $\sen x_{n}=0\to0$, mentre $\sen x_{n}'=1\to1$. Per il
coseno si usano le stesse successioni, scambiate di ruolo.
\end{esempio}

\begin{esempio}\label{es:sen-1-x}
Analogamente si dimostra la \eqref{eq:sen-1-x-anticipo}:
\begin{equation}\label{eq:sen-1-x}
  \lim_{x\to0}\sen\frac1x \qquad\text{non esiste.}
\end{equation}
La successione $x_{n}=\frac1{n\pi}$ converge a zero per $n\to+\infty$; in
corrispondenza la funzione $f(x)=\sen\frac1x$ assume i valori
\begin{equation}\label{eq:sen-1-x-a}
  f(x_{n})=\sen\frac1{x_{n}}=\sen(n\pi)=0 .
\end{equation}
Perciò $f(x_{n})$ è la successione costante uguale a zero, e converge a zero.
Analogamente, posto $x_{n}'=\frac1{\pi/2+2n\pi}$, risulta
\begin{equation}\label{eq:sen-1-x-b}
  f(x_{n}')=\sen\frac1{x_{n}'}=\sen\Bigl(\frac\pi2+2n\pi\Bigr)=1;
\end{equation}
quindi $x_{n}'\to0$, ma $f(x_{n}')\to1$, in contrasto con il fatto che $f(x_{n})\to0$.
Ne segue che $f(x)$ non ammette limite per $x\to0$ (si veda il grafico nella
\cref{fig:sen-1-x}). Con lo stesso argomento si vede che la funzione $g$ dell'avviso
«Una tabella non è una dimostrazione» del \cref{sec:premessa-funzioni} non ha limite per
$x\to0$: lungo $x_{m}=\frac1m$ si ha $\sen\frac\pi{x_{m}}=\sen(m\pi)=0$ e
$g(x_{m})=\frac{\sen x_{m}}{x_{m}}\to1$, mentre lungo $x_{m}'=\frac2{4m+1}$ si ha
$g(x_{m}')=\frac{\sen x_{m}'}{x_{m}'}+1\to2$.
\end{esempio}

Il programma seguente calcola i valori di $\sen\frac1x$ lungo le due successioni
dell'\cref{es:sen-1-x}.

\begin{python}[Due successioni, due limiti diversi]
import math

f = lambda x: math.sin(1 / x)

# x_n = 1/(n pi) e x'_n = 1/(pi/2 + 2 n pi) tendono entrambe a 0
print("      n      x_n        f(x_n)       x'_n        f(x'_n)")
for n in (1, 10, 100, 1000, 10**6):
    x = 1 / (n * math.pi)
    y = 1 / (math.pi / 2 + 2 * n * math.pi)
    print(f"{n:7d}   {x:.4e}  {f(x):+.2e}   {y:.4e}   {f(y):.12f}")
\end{python}
\begin{risultato}
      n      x_n        f(x_n)       x'_n        f(x'_n)
      1   3.1831e-01  +1.22e-16   1.2732e-01   1.000000000000
     10   3.1831e-02  -1.22e-15   1.5527e-02   1.000000000000
    100   3.1831e-03  +1.96e-15   1.5876e-03   1.000000000000
   1000   3.1831e-04  -3.21e-13   1.5912e-04   1.000000000000
1000000   3.1831e-07  -2.23e-10   1.5915e-07   1.000000000000
\end{risultato}

Si noti che il calcolatore non restituisce esattamente $\sen(n\pi)=0$: il numero $\pi$
è rappresentato con un errore relativo dell'ordine di $10^{-16}$, e l'argomento
$1/x_{n}\approx n\pi$ porta con sé un errore assoluto che cresce con $n$ (per
$n=10^{6}$ si arriva a $2\cdot10^{-10}$). È un'ulteriore conferma che la questione
dell'esistenza di un limite non si decide con il calcolo numerico: il programma
\emph{illustra} il ragionamento dell'\cref{es:sen-1-x}, ma non lo sostituisce.
